\documentclass[UTF8,12pt,fontset=fandol]{ctexart}

% 所有示例信息均为虚构内容，请替换为自己的信息。
\newcommand{\reporttitle}{通用工程课程报告模板}
\newcommand{\reportsubtitle}{General Engineering Report Template}
\newcommand{\coursename}{工程实践课程}
\newcommand{\departmentname}{示例工程学院}
\newcommand{\studentname}{张同学}
\newcommand{\studentid}{20260000000}
\newcommand{\classname}{示例一班}
\newcommand{\instructorname}{李老师}
\newcommand{\submissiondate}{\today}

\usepackage[a4paper,top=2.5cm,bottom=2.5cm,left=2.8cm,right=2.8cm]{geometry}
\usepackage{setspace}
\usepackage{fancyhdr}
\usepackage{graphicx}
\usepackage{booktabs}
\usepackage{array}
\usepackage{amsmath,amssymb}
\usepackage{siunitx}
\usepackage{xcolor}
\usepackage{listings}
\usepackage{tikz}
\usepackage{hyperref}

\hypersetup{
  unicode=true,
  colorlinks=true,
  linkcolor=blue!45!black,
  citecolor=green!35!black,
  urlcolor=blue!55!black,
  pdftitle={\reporttitle},
  pdfauthor={\studentname}
}

\setstretch{1.35}
\setlength{\parindent}{2em}
\setlength{\parskip}{0.25em}
\sisetup{detect-all}

\pagestyle{fancy}
\fancyhf{}
\fancyhead[L]{\coursename}
\fancyhead[R]{\studentname}
\fancyfoot[C]{\thepage}
\setlength{\headheight}{15pt}

\definecolor{codebackground}{RGB}{247,247,247}
\lstset{
  basicstyle=\ttfamily\small,
  backgroundcolor=\color{codebackground},
  frame=single,
  breaklines=true,
  numbers=left,
  numberstyle=\scriptsize,
  showstringspaces=false,
  tabsize=4
}

\newcommand{\metadatarow}[2]{%
  \makebox[4.2em][s]{#1} & \makebox[10cm][l]{\hrulefill\ #2\ \hrulefill}\\[1.2em]
}


\begin{document}

\begin{titlepage}
  \centering
  \vspace*{1.5cm}
  {\zihao{2}\bfseries \reporttitle\par}
  \vspace{0.5cm}
  {\Large \reportsubtitle\par}
  \vspace{2.5cm}
  \renewcommand{\arraystretch}{1.2}
  \begin{tabular}{rl}
    \metadatarow{课程名称}{\coursename}
    \metadatarow{学院}{\departmentname}
    \metadatarow{姓名}{\studentname}
    \metadatarow{学号}{\studentid}
    \metadatarow{班级}{\classname}
    \metadatarow{指导教师}{\instructorname}
    \metadatarow{提交日期}{\submissiondate}
  \end{tabular}
  \vfill
  {\small 本模板为非官方通用模板，请按学校或比赛要求调整。\par}
\end{titlepage}

\pagenumbering{Roman}
\section*{摘要}
本文展示一个可直接编译的通用工程课程报告结构，覆盖中文排版、公式、表格、流程图、代码和参考文献。使用者应把示例内容替换为自己的实验数据，并在提交前检查格式要求、来源和个人信息。

\noindent\textbf{关键词：}工程报告；\LaTeX；可复现性；实验记录

\section*{Abstract}
This document demonstrates a compilable engineering-report structure with equations, tables, diagrams, source code, and references. Replace all fictional metadata and sample values before submission.

\noindent\textbf{Keywords:} engineering report; \LaTeX; reproducibility; test record

\tableofcontents
\clearpage
\pagenumbering{arabic}

\section{项目背景与目标}
说明问题来源、约束条件、目标指标和验收方式。可量化目标应给出单位、容差和测试条件，而不是只写“性能良好”。

\section{方案设计}
\subsection{系统结构}
图~\ref{fig:flow} 使用 TikZ 直接绘制，因此模板不依赖外部示例图片。

\begin{figure}[htbp]
  \centering
  \begin{tikzpicture}[node distance=2.2cm,>=stealth,thick]
    \tikzstyle{block}=[draw,rounded corners,minimum width=2.8cm,minimum height=1cm,align=center]
    \node[block] (input) {输入与采样};
    \node[block,right of=input,xshift=1.2cm] (control) {算法与控制};
    \node[block,right of=control,xshift=1.2cm] (output) {执行与保护};
    \draw[->] (input) -- (control);
    \draw[->] (control) -- (output);
    \draw[->,bend left=35] (output.north) to node[above]{反馈} (input.north);
  \end{tikzpicture}
  \caption{通用闭环系统结构示例}
  \label{fig:flow}
\end{figure}

\subsection{理论模型}
一阶对象可写为
\begin{equation}
  G(s)=\frac{K}{\tau s+1},
  \label{eq:first-order}
\end{equation}
其中 $K$ 为稳态增益，$\tau$ 为时间常数。示例测量值可表示为 \SI{12.0}{\volt}，但实际报告必须标注仪器和测量不确定度。

\section{实现方法}
代码清单~\ref{lst:clamp} 展示一个具有明确边界的限幅函数。

\begin{lstlisting}[language=C,caption={限幅函数示例},label={lst:clamp}]
float clampf(float value, float minimum, float maximum) {
    if (value < minimum) return minimum;
    if (value > maximum) return maximum;
    return value;
}
\end{lstlisting}

\section{实验与结果}
表~\ref{tab:results} 中的数据仅用于展示排版，不能作为真实实验结论。

\begin{table}[htbp]
  \centering
  \caption{虚构测试数据示例}
  \label{tab:results}
  \begin{tabular}{cccc}
    \toprule
    工况 & 输入/V & 输出/V & 效率/\% \\
    \midrule
    A & 12.0 & 11.2 & 88.4 \\
    B & 15.0 & 14.1 & 90.1 \\
    C & 18.0 & 17.0 & 91.0 \\
    \bottomrule
  \end{tabular}
\end{table}

\section{讨论与局限}
说明误差来源、失败条件、尚未验证的假设和可能的改进。工程报告应保留失败数据与版本信息，以便他人复现。

\section{结论}
按目标逐项给出是否达成及证据位置，避免引入正文未验证的新结论。

\begin{thebibliography}{9}
\bibitem{lamport1994latex}
Leslie Lamport. \textit{LaTeX: A Document Preparation System}. Addison-Wesley, 2nd edition, 1994.
\end{thebibliography}

\appendix
\section{提交前检查}
\begin{itemize}
  \item 已替换姓名、学号、课程和教师等虚构信息；
  \item 图片、代码、数据和引用具有明确来源；
  \item 表格单位、有效数字和测试条件一致；
  \item PDF 中书签、目录、公式和参考文献均可正常显示。
\end{itemize}

\end{document}
